% GENERATED FILE. Edit canonical lessons, then run npm run generate:books.
% canonical-source-hash: fnv1a64:d2afc544acd6ed44
% canonical-lessons: PT-C311-infancia, PT-C311-adolescencia, PT-C311-juventude, PT-C311-velhice, PT-C311-geracao

\chapter{As idades da vida}
\label{ch:pt-as-idades-da-vida}
\hlchaptermodality{311}

\begin{chapteropening}
\textbf{By the end of this chapter:} \emph{I can talk about the ages of life: childhood, adolescence, youth, old age and generations.}

Complete the last lesson of chapter 311: I can talk about the ages of life: childhood, adolescence, youth, old age and generations.
\end{chapteropening}

\section[a infância]{a infância — childhood}

\label{lesson:PT-C311-infancia}



\begin{quote}
Before the new one: say the Portuguese for the godmother, then the Portuguese for to breastfeed.
\end{quote}

\subsection*{You'll want to know: a infância}
\textbf{a infância} — \textquotedblleft{}childhood\textquotedblright{}.

\begin{quote}
Passei a infância numa aldeia.
\end{quote}

\begin{quote}
\emph{I spent my childhood in a village.}
\end{quote}

\emph{Amigo de infância} — childhood friend.

\begin{cousinweb}
\emph{infância} is a learned borrowing from Latin \emph{īnfantia} \textquotedblleft{}inability to speak, childhood\textquotedblright{}, from \emph{īnfāns} \textquotedblleft{}not speaking\textquotedblright{}, built on \emph{fārī} \textquotedblleft{}to speak\textquotedblright{} (Lewis \& Short).
\end{cousinweb}

\begin{grammarlens}[title={What you've built}]
One more word for telling, explaining and describing at B1.
\end{grammarlens}

\subsection*{Your turn}
\begin{itemize}
\raggedright
  \item \emph{Say it:} \emph{a infância}
  \item \emph{Say it:} \emph{a infância}, once more
  \item \emph{Say it:} a sentence of your own with \emph{a infância}
  \item \emph{Recall:} say the Portuguese for the godmother, then the Portuguese for to breastfeed, then say \emph{a infância} again
\end{itemize}

\begin{tcolorbox}[breakable,colback=teal!4,colframe=teal!35!black,title={Before you move on}]
What does a infância mean? (\textquotedblleft{}Childhood\textquotedblright{}.) Say it once more.
\end{tcolorbox}
\section[a adolescência]{a adolescência — adolescence, the teenage years}

\label{lesson:PT-C311-adolescencia}



\begin{quote}
Before the new one: say the Portuguese for to breastfeed, then the Portuguese for childhood.
\end{quote}

\subsection*{You'll want to know: a adolescência}
\textbf{a adolescência} — \textquotedblleft{}adolescence, the teenage years\textquotedblright{}.

\begin{quote}
Na adolescência, eu era muito tímido.
\end{quote}

\begin{quote}
\emph{As a teenager I was very shy.}
\end{quote}

\emph{Um adolescente} — a teenager.

\begin{cousinweb}
\emph{adolescência} is a learned borrowing from Latin \emph{adolēscentia} \textquotedblleft{}youth\textquotedblright{}, from \emph{adolēscere} \textquotedblleft{}to grow up\textquotedblright{} (Lewis \& Short).
\end{cousinweb}

\begin{grammarlens}[title={What you've built}]
One more word for telling, explaining and describing at B1.
\end{grammarlens}

\subsection*{Your turn}
\begin{itemize}
\raggedright
  \item \emph{Say it:} \emph{a adolescência}
  \item \emph{Say it:} \emph{a adolescência}, once more
  \item \emph{Say it:} a sentence of your own with \emph{a adolescência}
  \item \emph{Recall:} say the Portuguese for to breastfeed, then the Portuguese for childhood, then say \emph{a adolescência} again
\end{itemize}

\begin{tcolorbox}[breakable,colback=teal!4,colframe=teal!35!black,title={Before you move on}]
What does a adolescência mean? (\textquotedblleft{}Adolescence, the teenage years\textquotedblright{}.) Say it once more.
\end{tcolorbox}
\section[a juventude]{a juventude — youth}

\label{lesson:PT-C311-juventude}



\begin{quote}
Before the new one: say the Portuguese for childhood, then the Portuguese for adolescence.
\end{quote}

\subsection*{You'll want to know: a juventude}
\textbf{a juventude} — \textquotedblleft{}youth\textquotedblright{}.

\begin{quote}
Na juventude, o meu avô foi pescador.
\end{quote}

\begin{quote}
\emph{In his youth my grandfather was a fisherman.}
\end{quote}

Also young people as a group: \emph{a juventude de hoje}.

\begin{cousinweb}
\emph{juventude} is a learned borrowing from Latin \emph{iuventūs, iuventūtis} \textquotedblleft{}youth\textquotedblright{}, from \emph{iuvenis} \textquotedblleft{}young\textquotedblright{} (Lewis \& Short).
\end{cousinweb}

\begin{grammarlens}[title={What you've built}]
One more word for telling, explaining and describing at B1.
\end{grammarlens}

\subsection*{Your turn}
\begin{itemize}
\raggedright
  \item \emph{Say it:} \emph{a juventude}
  \item \emph{Say it:} \emph{a juventude}, once more
  \item \emph{Say it:} a sentence of your own with \emph{a juventude}
  \item \emph{Recall:} say the Portuguese for childhood, then the Portuguese for adolescence, then say \emph{a juventude} again
\end{itemize}

\begin{tcolorbox}[breakable,colback=teal!4,colframe=teal!35!black,title={Before you move on}]
What does a juventude mean? (\textquotedblleft{}Youth\textquotedblright{}.) Say it once more.
\end{tcolorbox}
\section[a velhice]{a velhice — old age}

\label{lesson:PT-C311-velhice}



\begin{quote}
Before the new one: say the Portuguese for adolescence, then the Portuguese for youth.
\end{quote}

\subsection*{You'll want to know: a velhice}
\textbf{a velhice} — \textquotedblleft{}old age\textquotedblright{}.

\begin{quote}
A minha avó passou a velhice com a família.
\end{quote}

\begin{quote}
\emph{My grandmother spent her old age with the family.}
\end{quote}

Built on \emph{velho}. A gentler word is \emph{a terceira idade}.

\begin{cousinweb}
\emph{velhice} is built on \emph{velho} \textquotedblleft{}old\textquotedblright{}, which is inherited from Latin \emph{vetulus} \textquotedblleft{}a little old\textquotedblright{}, from \emph{vetus} \textquotedblleft{}old\textquotedblright{} (Lewis \& Short).
\end{cousinweb}

\begin{grammarlens}[title={What you've built}]
One more word for telling, explaining and describing at B1.
\end{grammarlens}

\subsection*{Your turn}
\begin{itemize}
\raggedright
  \item \emph{Say it:} \emph{a velhice}
  \item \emph{Say it:} \emph{a velhice}, once more
  \item \emph{Say it:} a sentence of your own with \emph{a velhice}
  \item \emph{Recall:} say the Portuguese for adolescence, then the Portuguese for youth, then say \emph{a velhice} again
\end{itemize}

\begin{tcolorbox}[breakable,colback=teal!4,colframe=teal!35!black,title={Before you move on}]
What does a velhice mean? (\textquotedblleft{}Old age\textquotedblright{}.) Say it once more.
\end{tcolorbox}
\section[a geração]{a geração — the generation}

\label{lesson:PT-C311-geracao}



\begin{quote}
Before the new one: say the Portuguese for youth, then the Portuguese for old age.
\end{quote}

\subsection*{You'll want to know: a geração}
\textbf{a geração} — \textquotedblleft{}the generation\textquotedblright{}.

\begin{quote}
Na minha geração, poucas pessoas tinham telemóvel.
\end{quote}

\begin{quote}
\emph{In my generation few people had a mobile.}
\end{quote}

\emph{De geração em geração} — from one generation to the next.

\begin{cousinweb}
\emph{geração} comes down from Old Portuguese \emph{geeraçon} (EtymDB), from Latin \emph{generātiō} \textquotedblleft{}a begetting\textquotedblright{}, built on \emph{genus} \textquotedblleft{}birth, kind\textquotedblright{}.
\end{cousinweb}

\begin{grammarlens}[title={What you've built}]
One more word for telling, explaining and describing at B1.
\end{grammarlens}

\subsection*{Your turn}
\begin{itemize}
\raggedright
  \item \emph{Say it:} \emph{a geração}
  \item \emph{Say it:} \emph{a geração}, once more
  \item \emph{Say it:} a sentence of your own with \emph{a geração}
  \item \emph{Recall:} say the Portuguese for youth, then the Portuguese for old age, then say \emph{a geração} again
\end{itemize}

\begin{tcolorbox}[breakable,colback=teal!4,colframe=teal!35!black,title={Before you move on}]
What does a geração mean? (\textquotedblleft{}The generation\textquotedblright{}.) Say it once more.
\end{tcolorbox}
